
\section{Adaptive Sparse Attention: Learned Per-Head Sparsity for Efficient Transformer Inference}


\textbf{Authors:}
\begin{itemize}
\item {{AUTHOR\_1\_NAME}}  -- {{AUTHOR\_1\_EMAIL}}
\end{itemize}
\begin{itemize}
\item {{AUTHOR\_2\_NAME}}  -- {{AUTHOR\_2\_EMAIL}}
\end{itemize}

\textbf{Affiliations:}
\begin{itemize}
\item   {{UNIVERSITY\_1}}, {{DEPARTMENT\_1}}
\end{itemize}
\begin{itemize}
\item   {{UNIVERSITY\_2}}, {{DEPARTMENT\_2}}
\end{itemize}

\textbf{Corresponding author:} {{CORRESPONDING\_AUTHOR}}
\textbf{Preprint date:} {{DATE}}
\textbf{Version:} v{{VERSION}}
\textbf{Code:} https://github.com/{{ORG}}/SVD\_Project/tree/main/framework/attentions/adaptive\_sparse\_attention.py
\textbf{Model weights:} https://huggingface.co/{{HF\_ORG}}
\textbf{License:} {{LICENSE}}



\subsection{Abstract}

Dense softmax attention computes $L^2$ pairwise scores per head, but empirical 
studies have repeatedly shown that the attention distribution is highly 
concentrated: only a small fraction of keys receive most of the attention mass. 
This suggests that most of the $L^2$ scores can be pruned with minimal impact 
on quality. Static sparsity patterns (fixed windows, fixed strides) fail to 
capture this variability, and existing dynamic methods either ignore per-head 
heterogeneity or introduce non-differentiable top-$k$ operations that hinder 
end-to-end training.

We introduce \textbf{Adaptive Sparse Attention} (ASA), a per-head learnable 
sparsification mechanism. Each of the $H$ heads maintains a single scalar 
logit $\ell_h$ controlling its retention budget: $k_h = \lfloor (1 - \sigma(\ell_h)) \cdot L \rfloor$. 
The top-$k$ selection is made differentiable via a \textbf{soft sigmoid threshold}: 
after computing the top-$k$-th score $\tau_h$, we apply a soft mask 
$\sigma((s_i - \tau_h) / T)$ that concentrates mass on the retained entries.

We evaluate ASA on DistilBERT, RoBERTa-base, and GPT-2 on SST-2. At default 
sparsity {{SPARSITY}} on decoders and {{SPARSITY\_ENC}} on encoders, ASA 
achieves \textbf{{{D\_ACC}}\%}, \textbf{{{R\_ACC}}\%}, and \textbf{{{G\_ACC}}\%} accuracy, within 
\textbf{{{D\_DELTA}}\%}, \textbf{{{R\_DELTA}}\%}, and \textbf{{{G\_DELTA}}\%} of the respective 
dense baselines. Attention FLOPs are reduced by up to \textbf{{{FLOP\_REDUCTION}}$\times$}. 
The learned per-head budgets show systematic structure: some heads converge 
to near-dense (sparsity $\approx$ 0) and others to near-sparse (sparsity $\approx$ 0.9), 
demonstrating head heterogeneity.

\textbf{Keywords:} sparse attention, learned sparsity, top-k, differentiable 
masking, per-head budgets



\subsection{Lay Summary}

Attention is the core operation in transformers. For a sequence of 1000 
words, it computes a million pairwise scores -- but in practice, most of these 
scores are near zero and contribute nothing. Throwing them away saves compute 
without hurting quality.

The challenge is \textit{which} scores to throw away. A fixed rule (e.g., "keep only 
nearby tokens") works poorly when important context is distant. Our method 
lets each attention head \textbf{learn} its own rule. A single trainable number 
per head controls what fraction of keys that head attends to. The learning 
is fully differentiable, so it can be trained end-to-end alongside the model.

We test this on three popular language models and find that the learned 
budgets are highly non-uniform: some heads remain dense, others become almost 
completely sparse. This shows that not all heads are equally important -- a 
fact that static sparsity methods ignore.



\subsection{Table of Contents}

1. [Introduction](01\_intro.md)
2. [Related Work](02\_related.md)
3. [Background](03\_background.md)
4. [Method](04\_method.md)
5. [Architecture Integration](05\_architecture.md)
6. [Experimental Setup](06\_setup.md)
7. [Results](07\_results.md)
8. [Analysis](08\_analysis.md)
9. [Discussion](09\_discussion.md)
10. [Limitations](10\_limitations.md)
11. [Conclusion](11\_conclusion.md)
12. [Reproducibility](12\_reproducibility.md)
13. [References](13\_references.md)

\textbf{Appendices:}
\begin{itemize}
\item [A. Mathematical Derivations](A\_math.md)
\end{itemize}
\begin{itemize}
\item [B. Implementation Details](B\_implementation.md)
\end{itemize}
\begin{itemize}
\item [C. Full Experimental Data](C\_data.md)
\end{itemize}
\begin{itemize}
\item [D. Negative Results](D\_negative\_results.md)
\end{itemize}
\begin{itemize}
\item [E. Ethical Considerations](E\_ethics.md)
\end{itemize}